\documentclass[10pt,a4paper]{article}
\usepackage[margin=25mm]{geometry}
\usepackage{amsmath,amssymb,amsthm}
\usepackage[colorlinks=true,urlcolor=blue]{hyperref}
\setlength{\emergencystretch}{3em}
\newtheorem{theorem}{Theorem}
\title{Four Rational Parking Functions, Not Twelve\\Counterexample to Conjecture 00000000456}
\author{idealistichacker\\\small AI-assisted proof and reproduction; no affiliation claimed}
\date{13 October 2026}
\begin{document}
\maketitle
\begin{abstract}
The stated factorial ratio is not the cardinality of standard rational parking
functions. In the height-first convention, parameters $(m,n)=(3,2)$ give four
labelled rational Dyck paths, while the displayed ratio is $4!/2!=12$.
An enumerative positive $q$-refinement cannot specialize at $q=1$ both to the
actual count and to this ratio. A self-contained pinned Lean project verifies
actual paths, every-vertex diagonal tests, exhaustive distinct labels, the
path/column bridge, the four-object count and the closed specialization negation.
\end{abstract}

\section{The named objects and the necessary clause}
Both source languages name rational parking functions and describe the factorial
ratio as their count. The Chinese explicitly states the counting formula; the
English calls the requested positivity a strengthening of that formula.
Following the height-first convention of \cite{ALW}, an $(m,n)$ rational parking
function is a Dyck path from $(0,0)$ to $(n,m)$, with its $m$ north steps labelled
bijectively by $1,\ldots,m$, increasing from bottom to top within each column.
We use positive coprime parameters $(m,n)=(3,2)$.

A genuine enumerative $q$-analogue $Q$ must satisfy $Q(1)=|\mathrm{PF}_{m,n}|$.
The stated counting formula would simultaneously require
\[
 Q(1)=\frac{(m+n-1)!}{(m-1)!}.
\]
This condition is independent of the chosen statistic. As a bare number, the
factorial ratio does have a positive $q$-integer product; that numerical fact
is not denied here. Such a product is not a refinement of the actual object
count in the counterexample. Dropping the object count would change the assertion.

\section{An exhaustive geometric count}
Let $c_i$ be the column where the $i$th north step starts, for $i=0,1,2$.
The path lies above the line $3x=2y$ precisely when
\[
 0\leq c_0\leq c_1\leq c_2<2,\qquad 3c_i\leq2i.
\]
Thus $(c_0,c_1,c_2)$ is either $(0,0,0)$ or $(0,0,1)$.
The column data reconstruct the unique word by inserting east runs before,
between and after the three north steps. Each horizontal run has its minimum
diagonal margin at the right endpoint; north steps increase the margin.
The terminal east run at height three ends at width two and stays above the
line. This accounts for every vertex and the endpoint, not just north steps.

For $(0,0,0)$ all three labels must increase, giving one labelling. For
$(0,0,1)$ only the first two must increase, giving three labellings:
\[
\begin{array}{c|c|c}
\text{word}&\text{north labels in path order}&\text{preferences in label order}\\\hline
NNNEE&(1,2,3)&(0,0,0)\\
NNENE&(1,2,3)&(0,0,1)\\
NNENE&(1,3,2)&(0,1,0)\\
NNENE&(2,3,1)&(1,0,0)
\end{array}
\]
Preferences record the north-step column for each label. Sorting columns and
sorting labels within each column reconstruct the labelled path; hence no
object or multiplicity is hidden by this representation.

\begin{theorem}
For standard rational parking functions at height-first parameters $(3,2)$,
the actual count is four and the stated factorial ratio is twelve.
Consequently the claimed enumerative positive $q$-refinement cannot exist.
\end{theorem}
\begin{proof}
The exhaustive count above gives four distinct objects, whereas
$(3+2-1)!/(3-1)!=24/2=12$. If a required refinement existed, its value at one
would be both four and twelve, a contradiction.
\end{proof}
The convention that lists width first also cannot rescue a universal formula:
the same paths at $(2,3)$ would give four against twenty-four. The Lean proof
advertises only the declared height-first witness, not an unverified extra case.

\section{Kernel-checked definition bridges}
The Lean source imports only the \texttt{Std} library bundled with Lean 4.33.1.
Actual north/east words have a recursive vertex test and exact step counts.
The theorem \texttt{word\_universe} covers every length-five word;
\texttt{raw\_word\_bridge} and \texttt{path\_injective} identify the path and
column representations without omitting terminal steps.
Every raw column/label object is enumerated, membership is equivalent to the
actual validity predicate, the four-element list is exact and has no duplicates.
Distinct labels in the three-element type use all labels; the same-column
strict inequalities are enforced before counting.

Finite natural-coefficient lists represent positive polynomials, with ordinary
Horner evaluation. The theorem \texttt{eval\_at\_one} identifies evaluation
with the coefficient sum. The final closed specialization theorem rules out
simultaneous evaluation to the genuine count and the proposed formula.
Kernel \texttt{decide} is used only on actual finite words, labels and their
fully specified predicates, not as a replacement for a definition bridge.
Every advertised theorem has its actual type and transitive axioms printed.

\section{Reproducibility and limits}
The accompanying Python reproducer independently checks all endpoint-compatible
words and label permutations, then cleanly rebuilds this pinned project and
audits every theorem. It has no external package dependencies or network
bootstrap requirement. Compiler binaries and bundled standard-library modules
remain trusted; no compiler audit, cross-platform run or human expert endorsement
is claimed. The submitted PDF is actually compiled from this source; validation
also recompiles and compares its bytes.

This is a counterexample to a necessary counting clause, not a denial of other
parking-function positivity theorems or of the bare numerical $q$-integer
product. No source conjecture or organizer metadata is modified, and no
affiliation, first-solution credit, score or upstream acceptance is asserted.

\begin{thebibliography}{1}
\bibitem{ALW} D. Armstrong, N. A. Loehr, G. S. Warrington,
\emph{Rational Parking Functions and Catalan Numbers},
arXiv:1403.1845v1 (2014), Sections 3.1--3.2 and Corollary 4.
\end{thebibliography}
\end{document}
